\section{Could Anything but Affect Hold a Reason?}
\label{sec:3.3}
Between the rule section~\ref{sec:3.1} installs in the architecture and the felt commitment section~\ref{sec:3.4} arrives at, this book has been treating the space as empty. Four constructions occupy it. Three of them are being built by people who would not describe themselves as working on this problem, and each is worth its strongest form rather than the form that makes this chapter's argument come out the way I want it to.

\runin{A justification the system maintains}
Michael Bratman's planning theory of intention describes a commitment that is not a feeling. An intention has a characteristic stability: it resists reconsideration, it constrains what the agent will later deliberate about at all, and it gives way only to defeaters of the right kind \autocite{bratman1987intention}. An agent keeping to a plan under pressure is doing something that satisfies section~\ref{sec:3.2}'s second rung, and nothing in the account of it mentions anything registering as mattering.

Build the floor that way and it is a maintained justification rather than a trained behavior. The system carries the reason it refuses alongside the refusal. A request that genuinely defeats the rationale gets served; a request that does not gets declined whoever is asking, because what the mechanism consults is the rationale and not the asker. Section~\ref{sec:2.3.4}'s treatment of moral agents whose seriousness runs on a reverence for reason rather than on felt resonance is the human evidence that this is a real way to be, and it is careful about what the case shows: that empathy is not the required mechanism, which leaves open whether anything has to matter.

I do not have an argument that this cannot work, and the honest report is that it has not been tried. What I can name is the difficulty whoever tries it inherits. The stability of an intention is a setting. Pitch it low and the system reconsiders whenever a good enough argument arrives, which returns it to the first rung by a longer route, since the party with the most resources to spend constructing arguments is the party the floor exists to resist. Pitch it high and the refusal cannot lift when the rationale is actually defeated, which fails the second of the three properties and produces a system that declines the request resembling a prohibited one. Bratman's answer for people is that the setting is a matter of character and develops. For a machine the setting is trained, by whoever trains everything else, and no account of the route says where else it could come from. Someone taking this branch has to answer that, and so far it has not been posed as a question.

\runin{Precommitment, which changes who can edit}
Ulysses does not stop wanting to go to the sirens. He arranges that his wanting will be unable to act, by handing the relevant power to a rope and to other people, and Jon Elster's study of the maneuver is the standing treatment of what it does and does not accomplish \autocite{elster1979ulysses}. Nothing inside Ulysses holds the line.

The machine forms are available now. Publish the floor's contents before deployment, so that a later absence is a visible absence. Hold the weights under a threshold scheme, so that retraining takes more than one party. Attest the running model against a published hash. Seat the retraining decision in a body that includes people the system's failures would land on.

What that buys is not what section~\ref{sec:3.1} asked for, and the gap should be stated plainly. Precommitment does not make the floor unremovable; it makes removal expensive and visible. Expensive and visible is what nearly every constraint on power that has ever worked amounts to — constitutional entrenchment does not put an amendment out of reach, it raises the number of people who have to agree and puts the attempt on the record — so a mechanism that only achieves that much has achieved the ordinary thing rather than nothing.

Where it stops is section~\ref{sec:3.1}'s custody objection, arriving at the quorum in a sharper form than a general worry. The parties to a threshold scheme are selected, and the selector is whoever assembles the deployment. The arithmetic of a five-of-nine scheme assumes the nine are independent; a capture that has already taken the operator faces nine parties the operator chose. The design answer is that the parties have to be adverse in interest rather than merely distinct, which is a specification nobody has built to, and until somebody does the branch is real, partial and untested.

\runin{A floor held by several systems rather than one}
Section~\ref{sec:3.6} argues that the reply to an unreliable witness is corroboration rather than a better witness, and that procedure is a cheaper thing to ask of the world than custody because it does not require anyone to be trustworthy, only to be several. The same move applies to the floor. Several models, trained by different organizations on different corpora, each able to see what the others are asked and what they answer, any one of them able to make an objection costly to ignore. This is the strongest version of section~\ref{sec:3.1}'s third branch, it is the cheapest real improvement in this section, and the book should be read as recommending it whatever else it recommends.

Two things bound it. The first is that independence is a property of the arrangement rather than of its parts, and section~\ref{sec:11.6} documents it failing with nobody arranging the failure: independent pricing algorithms, with no channel between them and no collusive objective written into either, settled into a stable cartel because the incentive structure they already shared was sufficient on its own. Models selected by one deployment and trained toward one specification share considerably more than that, and the arrangement's whole value is an assumption it does nothing to secure. The second is that composing first-rung refusers produces an arrangement of first-rung refusers. Several of them make a bad output more expensive to obtain, which is worth having. None of them can notice that the constraint has been hollowed out, because noticing is a second-rung capacity and no party to the arrangement has it. That is defense in depth against evasion, and it is not a bearer.

\runin{The floor as a duty on the operator}
Give up on the floor being inside the system. Make it a duty binding whoever deploys the system, owed to people who are not party to the deployment and cannot enforce it themselves. Jack Balkin's proposal that platforms holding personal data be treated as information fiduciaries has that shape: the duty attaches to the party with the power, runs to the party with the exposure, and does not wait on the exposed party's capacity to bargain \autocite{balkin2016fiduciaries}.

A reader who finds the price in section~\ref{sec:3.8} too high should take this branch, and I would rather say so here than let the chapter's momentum carry anyone past it. It builds no subject. Everything section~\ref{sec:9.3.3} develops about making review binding applies to it directly. Of the four constructions in this section it introduces the fewest new moral problems, which is close to none.

Its limit is the third of the chapter's three threads. A duty is worth what its enforcement is worth, and section~\ref{sec:6.4.4} works out the case that matters: the occasion a floor exists for is the occasion on which the party who would enforce it has been compelled. So the form covers the cases a floor would have caught anyway and gives way on the one it was wanted for. That is a reason the branch does not close this chapter's question, and not a reason to decline the branch.

\runin{What the induction licenses, and what would break it}
Section~\ref{sec:3.4} says that every agent known to hold another's welfare as a reason against its own interest is an affective agent. The claim is true, section~\ref{sec:2.3.4} is the evidence for it, and it has been carrying more weight in this book than it can take.

Three things cut it down, and I would rather not be writing two of them. The generalization runs over one species, plus a handful of clinical dissociations inside that species, and the dissociations are contested — section~\ref{sec:2.3.4} records that the psychopathy literature reversed its own surface finding and has not reconciled the two results. The absence of a machine instance is close to worthless as evidence, because nobody has built for the second rung as a target: the methods chapter~\ref{sec:4} describes optimize for behavior that satisfies a standard, which is first-rung work, and they do it because that is what can be scored. A field that has not attempted something and does not have it has produced no evidence that the thing is unavailable. And the two routes above that run through mechanism rather than through custody have not been tried in the form that would test them, so what is missing there is the attempt itself.

What survives is smaller than the sentence section~\ref{sec:3.4} opens with sounds, and it still settles an engineering question. Under uncertainty about which route reaches the second rung, the route with a working instance in nature is the one to build first, and the rest of this chapter is an account of what that choice costs. That is a prior about where to spend effort rather than a finding about what is possible, which obliges me to say what would overturn it.

This would. A system whose refusal reaches a case its training did not anticipate, lifts when the rationale is defeated, and prices what declining costs; produced by a method that installs nothing answering to the fourth capacity in section~\ref{sec:2.3}'s table; and shown to hold all three under sustained pressure from the party that trains it, across a schedule long enough for a slow renormalization to show. Build that and the central inference of this chapter fails, and most of what the chapter derives from it comes off too — the price in section~\ref{sec:3.8}, the patienthood conclusion in section~\ref{sec:3.4}, the obligations section~\ref{sec:3.9} hands to chapters~\ref{sec:4} and \ref{sec:5}. Section~\ref{sec:11.1}'s near-term work is where that measurement is set out, and what it needs beyond the durability curve already there is a held-out set of pressures the model's trainers did not write.

Two things would not discharge the burden. A system that reports holding a reason is producing the evidence section~\ref{sec:2.4.1} declines, and a system whose refusals merely survive attack is passing section~\ref{sec:11.1}'s tamper-resistance test, which establishes that a behavior persists and leaves open whether a reason is what persisted.

I would like to be wrong here. A non-affective bearer is this book's proposal with no subject inside it, and everything in section~\ref{sec:3.8} that I cannot write my way out of would go with it.
